\documentclass[11pt]{article}

\usepackage[spanish]{babel}
\usepackage[utf8]{inputenc}
\usepackage[T1]{fontenc}
\usepackage{amsmath, amssymb, amsthm, mathtools}
\usepackage{geometry}
\usepackage{xcolor}
\usepackage{tcolorbox}

\geometry{letterpaper, margin=2.2cm}

% ==============================================================================
% DEFINICIONES DE ENTORNOS PERSONALIZADOS PARA COMPILACIÓN
% (Permite compilar el archivo a PDF sin alterar la estructura que lee la IA)
% ==============================================================================

\newenvironment{motivacion}{%
  \section*{1. Motivación}%
}{}

\newenvironment{preguntaguia}{%
  \begin{tcolorbox}[colback=cyan!5,colframe=cyan!75!black,title=Pregunta Guía]%
}{%
  \end{tcolorbox}%
}

\newenvironment{teoria}{%
  \section*{2. Definición y Teoría}%
}{}

\newenvironment{definicion}[1]{%
  \begin{tcolorbox}[colback=blue!5,colframe=blue!75!black,title=Definición: #1]%
}{%
  \end{tcolorbox}%
}

\newenvironment{teorema}[1]{%
  \begin{tcolorbox}[colback=green!5,colframe=green!75!black,title=Teorema: #1]%
}{%
  \end{tcolorbox}%
}

\newenvironment{alerta}[1]{%
  \begin{tcolorbox}[colback=yellow!10,colframe=orange!80!black,title=Alerta: #1]%
}{%
  \end{tcolorbox}%
}

\newenvironment{procesamiento}[1]{%
  \begin{tcolorbox}[colback=gray!5,colframe=gray!60!black,title=Procedimiento: #1]%
}{%
  \end{tcolorbox}%
}

\newenvironment{aplicacion}{%
  \section*{3. Aplicación y Práctica}%
}{}

\newenvironment{ejemplo}[1]{%
  \begin{tcolorbox}[colback=white,colframe=gray!40,title=Ejemplo: #1]%
}{%
  \end{tcolorbox}%
}

% Comandos para preguntas interactivas

\newenvironment{preguntaalternativas}[1]{\subsection*{Pregunta: #1}\par\vspace{0.1cm}}{\vspace{0.3cm}}
\newcommand{\opcion}[3]{\par\noindent $\bullet$ \textbf{#1} \textit{(#2)} - #3\par\vspace{0.15cm}}

\newenvironment{preguntaverdaderofalso}[1]{\subsection*{Pregunta V/F: #1}\par\vspace{0.1cm}}{\vspace{0.3cm}}
\newcommand{\verdaderofalso}[3]{\par\noindent\textbf{Respuesta correcta: #1} \\ \textit{Si marca Falso:} #2 \\ \textit{Si marca Verdadero:} #3\par\vspace{0.15cm}}

\newenvironment{preguntacasillas}[1]{\subsection*{Selección Múltiple: #1}\par\vspace{0.1cm}}{\vspace{0.3cm}}
\newcommand{\casilla}[3]{\par\noindent $\square$ \textbf{#1} \textit{(#2)} - #3\par\vspace{0.15cm}}

% --- TÉRMINOS PAREADOS (2 COLUMNAS \colon NÚMEROS Y LETRAS) ---
\newenvironment{pareadosdoscolumnas}[1]{%
  \begin{tcolorbox}[colback=violet!5,colframe=violet!75!black,title=Términos Pareados (2 Columnas): #1]%
}{%
  \end{tcolorbox}%
}


% Entorno Términos Pareados (3 Columnas)
\newenvironment{pareadostrescolumnas}[1]{\subsection*{Términos Pareados: #1}}{\vspace{0.3cm}}
\newcommand{\columnaI}[1]{\textbf{Columna 1:}\begin{enumerate}#1\end{enumerate}}
\newcommand{\columnaII}[1]{\textbf{Columna 2:}\begin{itemize}#1\end{itemize}}
\newcommand{\columnaIII}[1]{\textbf{Columna 3:}\begin{itemize}#1\end{itemize}}
\newcommand{\pareotres}[2]{\textbf{Cruce #1:} #2\par}

% Entornos para ejercicios
\newenvironment{ejercicios}{%
  \section*{4. Ejercicios Resueltos y Propuestos}%
}{}

\newenvironment{ejercicioresuelto}[2]{\subsection*{Ejercicio Resuelto: #1 \small{[#2]}}}{}
\newenvironment{ejerciciodemostracion}[2]{\subsection*{Demostración: #1 \small{[#2]}}}{}
\newenvironment{ejerciciopropuesto}[2]{\subsection*{Ejercicio Propuesto: #1 \small{[#2]}}}{}

\newcommand{\enunciado}[1]{\textbf{Enunciado:} #1\par\vspace{0.2cm}}
\newcommand{\solucion}[1]{\textbf{Solución:} #1\par\vspace{0.4cm}}
\newcommand{\demostracion}[1]{\textbf{Demostración:} #1\par\vspace{0.4cm}}
\newcommand{\pista}[1]{\textbf{Pista:} #1\par\vspace{0.4cm}}

% Entorno Fórmulas
\newenvironment{formulas}{\section*{5. Fórmulas Clave}\begin{description}}{\end{description}}
\newcommand{\formula}[3]{\item[#1:] $#2$ \\ \textit{#3} \vspace{0.2cm}}

% ==============================================================================
% 1. METADATOS TÉCNICOS DE DESPLIEGUE (COMPLETAR OBLIGATORIAMENTE)
% ==============================================================================
% ID_CURSO \colon calculo-multivariable
% NUMERO_UNIDAD: 1
% TITULO_UNIDAD \colon Funciones de Varias Variables
% NUMERO_CAPITULO: 1.1
% TITULO_CAPITULO \colon Campos Escalares
% ESTADO_BLOQUEADO \colon false
% ESTADO_COMPLETADO \colon false
% ==============================================================================

\begin{document}

% ==============================================================================
% PESTAÑA 1 \colon MOTIVACIÓN (contentMotivation)
% ==============================================================================
\begin{motivacion}
  \textbf{¿Cómo modelamos el espacio?}
  
  Imagina que estás en tu habitación en este instante. Si quisiéramos registrar la temperatura en cada rincón, no nos basta con un único número. La temperatura varía si te acercas a la ventana, si subes al techo o si te sientas en el suelo.
  
  Para describir esto matemáticamente, debemos asignar a cada coordenada espacial $(x, y, z)$ un único valor térmico $T$. Así, la temperatura es una función $T(x,y,z)$. 
  
  \textbf{Concepto clave:} Un campo escalar es simplemente una regla que le asigna un único número real (temperatura, presión, altitud, densidad) a cada punto del espacio.

  \begin{preguntaguia}
    ¿Cómo cambia el cálculo cuando una función depende de más de una variable?
  \end{preguntaguia}
\end{motivacion}


% ==============================================================================
% PESTAÑA 2 \colon DEFINICIÓN Y TEORÍA (contentTheory)
% ==============================================================================
\begin{teoria}
  Para transitar del cálculo en una variable al cálculo multivariable, debemos formalizar matemáticamente la noción de asignar valores numéricos a puntos en el espacio. A continuación, estableceremos los pilares algebraicos de los campos escalares.

  % --- DEFINICIÓN 1 : CAMPO ESCALAR ---
  \begin{definicion}{Campo Escalar}
    Un \textbf{campo escalar} (o función de varias variables) es una función $f$ que asigna a cada elemento o punto $\vec{x} = (x_1, x_2, \dots, x_n)$ de un subconjunto $D \subseteq \mathbb{R}^n$ un único número real denotado por $f(\vec{x})$ o $f(x_1, x_2, \dots, x_n)$.
    
    Matemáticamente lo expresamos como:
    $$ f \colon D \subseteq \mathbb{R}^n \to \mathbb{R} $$
    Donde $n$ representa la dimensión del espacio de entrada (habitualmente $n=2$ o $n=3$).
  \end{definicion}

  % --- DEFINICIÓN 2 : DOMINIO ---
  \begin{definicion}{Dominio de un Campo Escalar}
    El \textbf{dominio} de un campo escalar $f$, denotado por $\operatorname{dom}(f)$, es el conjunto de todos los puntos en $\mathbb{R}^n$ para los cuales la regla de correspondencia de la función produce un número real bien definido. 
    
    Formalmente se define como:
    $$ \operatorname{dom}(f) = \{ \vec{x} \in \mathbb{R}^n \colon f(\vec{x}) \in \mathbb{R} \} $$
    Si el dominio no se especifica explícitamente, se asume que es el \textit{dominio natural}, es decir, el conjunto más grande posible de puntos donde la expresión matemática es válida.
  \end{definicion}

  % --- ALERTA DE DOMINIO ---
  \begin{alerta}{Visualización del Dominio}
    ¡Cuidado con la geometría! A diferencia de Cálculo I donde el dominio es un intervalo o unión de intervalos en la recta real ($\mathbb{R}$), en Cálculo Multivariable el dominio es una \textbf{región geométrica} en el plano XY ($\mathbb{R}^2$), en el espacio XYZ ($\mathbb{R}^3$) o en hiperespacios de mayor dimensión.
  \end{alerta}

  % --- DEFINICIÓN 3 : IMAGEN ---
  \begin{definicion}{Imagen o Recorrido}
    La \textbf{imagen} (o recorrido) de un campo escalar $f$, denotada por $\operatorname{im}(f)$ o $\text{Rec}(f)$, es el conjunto de todos los valores reales que la función efectivamente toma a medida que evaluamos todos los puntos pertenecientes al dominio.
    
    Formalmente se define como:
    $$ \operatorname{im}(f) = \{ y \in \mathbb{R} \colon y = f(\vec{x}) \text{ para algún } \vec{x} \in \operatorname{dom}(f) \} $$
  \end{definicion}

  % --- PROCEDIMIENTO / TRAMPA COGNITIVA ---
  \begin{procesamiento}{Restricciones Clásicas de Dominio}
    Para determinar analíticamente el dominio natural de un campo escalar, debes buscar y plantear las mismas restricciones algebraicas que conoces de funciones reales \colon \begin{itemize}
      \item Denominadores \colon Deben ser distintos de cero ($g(\vec{x}) \neq 0$).
      \item Raíces de índice par \colon El radicando debe ser no negativo ($g(\vec{x}) \geq 0$).
      \item Logaritmos \colon El argumento debe ser estrictamente positivo ($g(\vec{x}) > 0$).
    \end{itemize}
    La resolución de estas inecuaciones vectoriales definirá las fronteras de tu región matemática en el espacio.
  \end{procesamiento}
\end{teoria}


% ==============================================================================
% PESTAÑA 3 \colon APLICACIÓN Y PRÁCTICA (contentApplication)
% ==============================================================================
\begin{aplicacion}
  Ahora que entendemos qué es un campo escalar y su dominio, veamos cómo se comporta esto en la práctica con una función de dos variables y pongamos a prueba tus habilidades de análisis geométrico.

  % --- EJEMPLO SIMPLE ---
  \begin{ejemplo}{Evaluando un Campo Escalar y su Dominio}
    Consideremos el campo escalar dado por la función:
    $$ f(x,y) = \sqrt{9 - x^2 - y^2} $$
    
    \textbf{1. Evaluando un punto:}
    Si evaluamos la función en las coordenadas $(1, 2)$, obtenemos:
    $$ f(1,2) = \sqrt{9 - (1)^2 - (2)^2} = \sqrt{9 - 1 - 4} = \sqrt{4} = 2 $$
    Esto significa que en el punto $(1,2)$ del plano, nuestro campo escalar (que podría representar, por ejemplo, la altura de una superficie o domo hemisférico) tiene un valor de 2.
    
    \textbf{2. Determinando el Dominio natural:}
    Siguiendo nuestra trampa cognitiva sobre las restricciones, sabemos que el interior de una raíz cuadrada de índice par debe ser mayor o igual a cero:
    $$ 9 - x^2 - y^2 \geq 0 \implies x^2 + y^2 \leq 9 $$
    Por lo tanto, el dominio de $f$ es el conjunto de todos los puntos $(x,y)$ geométricamente ubicados dentro y sobre la frontera de una circunferencia de radio 3 centrada en el origen. ¡Cualquier punto fuera de ese círculo anulará nuestra función en los números reales!
  \end{ejemplo}

  % --- TIPO 1 \colon PREGUNTA DE ALTERNATIVAS (Cálculo analítico de dominio) ---
  \begin{preguntaalternativas}{Dominio con Logaritmos}
    ¿Cuál es el dominio natural del campo escalar $f(x,y) = \ln(y - x^2)$?
    \opcion{Todos los puntos $(x,y)$ tales que $y > x^2$.}{Correcto}{¡Excelente! El argumento de un logaritmo natural debe ser estrictamente positivo. Geométricamente, esto representa todos los puntos estrictamente por encima de la parábola $y = x^2$.}
    \opcion{Todos los puntos $(x,y)$ tales que $y \geq x^2$.}{Incorrecto}{¡Cuidado! El logaritmo natural no está definido para el valor cero. La desigualdad debe ser estricta.}
    \opcion{Todos los puntos $(x,y)$ tales que $y \neq x^2$.}{Incorrecto}{Esa sería la restricción si el término estuviera en un denominador. Recuerda que un logaritmo tampoco admite valores negativos en su argumento.}
  \end{preguntaalternativas}

  % --- TIPO 2 \colon PREGUNTA DE VERDADERO O FALSO (Concepto teórico) ---
  \begin{preguntaverdaderofalso}{Naturaleza del Campo Escalar}
    Un campo escalar definido como $f \colon \mathbb{R}^3 \to \mathbb{R}$ toma una coordenada del espacio tridimensional como entrada y produce un vector como salida.
    
\verdaderofalso{F}{Incorrecto. Recuerda la palabra 'escalar'. La salida es siempre un único número real (como una temperatura o una altitud), no un vector.}{¡Correcto! Entregamos un punto o vector de entrada $(x,y,z)$, pero la función nos devuelve un escalar (un único número real).}
  \end{preguntaverdaderofalso}

  % --- TIPO 3 \colon PREGUNTA DE CASILLAS (Combinación de restricciones) ---
  \begin{preguntacasillas}{Múltiples Restricciones}
    Considera la función $h(x,y) = \frac{\sqrt{x+y}}{x-1}$. Selecciona \textbf{todas} las condiciones matemáticas que deben cumplirse simultáneamente para que un punto $(x,y)$ pertenezca a su dominio natural \colon \casilla{$x + y \geq 0$}{Correcto}{Correcto. Al ser una raíz cuadrada en el numerador, su interior no puede ser negativo, pero sí puede ser cero.}
    \casilla{$x \neq 1$}{Correcto}{Correcto. El denominador completo debe ser distinto de cero para evitar la división por cero.}
    \casilla{$x + y > 0$}{Incorrecto}{Incorrecto. Si bien el interior no puede ser negativo, en este caso sí está permitido que sea cero porque la raíz está en el numerador.}
    \casilla{$y \neq 1$}{Incorrecto}{Incorrecto. La variable $y$ puede tomar el valor 1 sin problemas, siempre y cuando no se infrinja la condición de la raíz.}
  \end{preguntacasillas}




  % --- TIPO 4 \colon TÉRMINOS PAREADOS (3 Columnas - 5 Opciones) ---
  \begin{pareadostrescolumnas}{Asociación Avanzada \colon Función, Dominio e Imagen}
    
    % COLUMNA 1 : Funciones (La IA mapeará Números 1 al 5)
    \columnaI{
      \item $f(x,y) = \sqrt{1 - x^2 - y^2}$
      \item $f(x,y) = \ln(x + y)$
      \item $f(x,y) = e^{-x^2 - y^2}$
      \item $f(x,y) = \frac{1}{x^2 + y^2}$
      \item $f(x,y) = \sqrt{x^2 + y^2 - 1}$
    }
    
    % COLUMNA 2 : Dominios (La IA mapeará Letras A a E)
    \columnaII{
      \item $D = \{(x,y) \in \mathbb{R}^2 \colon x^2 + y^2 \geq 1\}$
      \item $D = \mathbb{R}^2 \setminus \{(0,0)\}$
      \item $D = \{(x,y) \in \mathbb{R}^2 \colon y > -x\}$
      \item $D = \{(x,y) \in \mathbb{R}^2 \colon x^2 + y^2 \leq 1\}$
      \item $D = \mathbb{R}^2$
    }
    
    % COLUMNA 3 : Imágenes/Recorridos (La IA mapeará Romanos I a V)
    \columnaIII{
      \item $\operatorname{im}(f) = \mathbb{R}$
      \item $\operatorname{im}(f) = [0, 1]$
      \item $\operatorname{im}(f) = [0, \infty)$
      \item $\operatorname{im}(f) = (0, 1]$
      \item $\operatorname{im}(f) = (0, \infty)$
    }

    % --- SOLUCIONES Y RETROALIMENTACIÓN ---
    \pareotres{1-D-II}{¡Correcto! La raíz exige $1 - (x^2+y^2) \geq 0$, un disco cerrado de radio 1. Su máximo es 1 (en el origen) y su mínimo es 0 (en la frontera).}
    \pareotres{2-C-I}{¡Excelente! El argumento $x+y>0$ define un semiplano abierto. Al no estar acotado, el logaritmo recorre todos los números reales.}
    \pareotres{3-E-IV}{¡Muy bien! No hay restricciones en el exponente, el dominio es todo $\mathbb{R}^2$. Su máximo valor es $e^0 = 1$, y decae hacia 0 sin tocarlo.}
    \pareotres{4-B-V}{¡Perfecto! El único problema ocurre si el denominador es cero, por lo que excluimos solo el origen $(0,0)$. Como la suma de cuadrados siempre es positiva, la fracción crece hacia infinito y nunca toma el valor cero ni negativos.}
    \pareotres{5-A-III}{¡Correcto! A diferencia de la función 1, aquí necesitamos que $x^2 + y^2 - 1 \geq 0$, lo que geométricamente es el exterior de un disco de radio 1 (incluyendo la frontera). Como las variables pueden crecer indefinidamente, la raíz va desde 0 hasta infinito.}

  \end{pareadostrescolumnas}

\end{aplicacion}


% ==============================================================================
% PESTAÑA 4 \colon EJERCICIOS RESUELTOS Y PROPUESTOS (contentExercises)
% ==============================================================================
\begin{ejercicios}

  % --- 1. EJERCICIO RESUELTO (Analítico / Intermedio) ---
  \begin{ejercicioresuelto}{Determinación de Dominio con Múltiples Restricciones}{nivel-2}
    \enunciado{
      Determine analíticamente y describa geométricamente el dominio natural del campo escalar dado por la expresión:
      $$f(x,y) = \frac{\sqrt{9 - x^2 - y^2}}{\ln(y - x)}$$
    }
    \solucion{
      Para que la función entregue un valor real bien definido, debemos plantear y resolver simultáneamente tres restricciones algebraicas \colon \textbf{Restricción 1 (Raíz cuadrada):} El radicando del numerador no puede ser negativo:
      $$ 9 - x^2 - y^2 \geq 0 \implies x^2 + y^2 \leq 9 $$
      Geométricamente, esto representa un disco cerrado de radio $3$ centrado en el origen $(0,0)$.
      
      \textbf{Restricción 2 (Logaritmo natural):} El argumento del logaritmo en el denominador debe ser estrictamente positivo:
      $$ y - x > 0 \implies y > x $$
      Esto corresponde al semiplano abierto ubicado estrictamente por encima de la recta identidad $y = x$.
      
      \textbf{Restricción 3 (Denominador no nulo):} El denominador completo no puede anularse:
      $$ \ln(y - x) \neq 0 \implies y - x \neq e^0 \implies y - x \neq 1 \implies y \neq x + 1 $$
      Esto significa que debemos excluir todos los puntos que pertenecen a la recta transladada $y = x + 1$.
      
      \textbf{Conclusión y Descripción del Dominio:}
      El dominio natural del campo escalar es la intersección de estas tres regiones:
      $$ \operatorname{dom}(f) = \left\{ (x,y) \in \mathbb{R}^2 \;\colon\; x^2 + y^2 \leq 9 \;\land\; y > x \;\land\; y \neq x + 1 \right\} $$
      Geométricamente, corresponde a la mitad superior-izquierda del círculo de radio 3 (cortado por la recta $y=x$), excluyendo los puntos de la frontera sobre dicha recta y quitando completamente el segmento de la recta $y = x + 1$ que cruza por dentro de la figura.
    }
  \end{ejercicioresuelto}

  % --- 2. EJERCICIO DE DEMOSTRACIÓN (Teórico / Avanzado) ---
  \begin{ejerciciodemostracion}{Invarianza por Simetría Radial}{nivel-3}
    \enunciado{
      Un campo escalar $f \colon \mathbb{R}^2 \to \mathbb{R}$ posee \textit{simetría radial} si su valor depende únicamente de la distancia del punto al origen. Es decir, si existe una función de una variable $g \colon [0, \infty) \to \mathbb{R}$ tal que $f(x,y) = g(\sqrt{x^2+y^2})$. 
      
      Demuestre rigurosamente que el campo escalar $f(x,y) = \ln(1 + x^2 + y^2)$ posee simetría radial, determine explícitamente la función $g(t)$ asociada, y pruebe analíticamente que la imagen del campo es $\operatorname{im}(f) = [0, \infty)$.
    }
    \demostracion{
      \textbf{Parte 1 \colon Demostración de Simetría Radial}
      Definamos la variable $t = \sqrt{x^2 + y^2}$, la cual representa la distancia euclidiana de cualquier punto $(x,y)$ al origen $(0,0)$. Dado que las variables están en los números reales, al elevar al cuadrado obtenemos $t^2 = x^2 + y^2$. 
      
      Sustituyendo directamente en la regla de correspondencia de nuestro campo escalar:
      $$ f(x,y) = \ln(1 + (x^2 + y^2)) = \ln(1 + t^2) $$
      Como la expresión resultante depende única y exclusivamente del parámetro de distancia $t$, queda demostrado que $f$ posee simetría radial. La función unidimensional asociada es:
      $$ g(t) = \ln(1 + t^2) \quad \text{con } t \in [0, \infty) $$
      
      \textbf{Parte 2 \colon Determinación Rigurosa de la Imagen}
      Para hallar la imagen, analizamos el comportamiento de $g(t)$ en su dominio restringido $[0, \infty)$:
      \begin{enumerate}
        \item Dado que $t \geq 0$, entonces $t^2 \geq 0$, lo que implica que $1 + t^2 \geq 1$.
        \item Aplicando la función logaritmo natural (que es estrictamente creciente en todo su dominio) a la desigualdad anterior, obtenemos:
        $$ \ln(1 + t^2) \geq \ln(1) \implies g(t) \geq 0 $$
        \item El valor mínimo absoluto es $0$ y se alcanza únicamente en $t=0$ (es decir, en el origen $f(0,0) = 0$).
        \item Como $\lim_{t \to \infty} \ln(1 + t^2) = \infty$ y la función es continua, por el Teorema del Valor Intermedio el recorrido toma todos los valores intermedios.
      \end{enumerate}
      Por lo tanto, la imagen del campo escalar es, formalmente, $\operatorname{im}(f) = [0, \infty)$.
    }
  \end{ejerciciodemostracion}

  % --- 3. EJERCICIO PROPUESTO 1 \colon Hipérbola Abierta ---
  \begin{ejerciciopropuesto}{Restricciones Hiperbólicas en el Plano}{nivel-2}
    \enunciado{
      Considere el campo escalar definido por la regla de correspondencia:
      $$ h(x,y) = \ln(x \cdot y - 1) $$
      Determine su dominio natural y describa las fronteras que delimitan esta región matemática en el plano cartesiano.
    }
    \pista{
      \textbf{Pista metodológica:} La restricción del logaritmo exige que el producto de las variables cumpla $x \cdot y > 1$. Analiza este comportamiento separando el análisis para cuando $x > 0$ y cuando $x < 0$. Recuerda que la frontera matemática está dada por las dos ramas de la hipérbola equilátera $y = \frac{1}{x}$, y que el dominio consta de dos regiones disjuntas en el primer y tercer cuadrante.
    }
  \end{ejerciciopropuesto}

  % --- 4. EJERCICIO PROPUESTO 2 \colon Logaritmo Cono Abierto ---
  \begin{ejerciciopropuesto}{Restricción Logarítmica e Hiperbólica}{nivel-2}
    \enunciado{
      Determine analíticamente el dominio natural y la imagen del campo escalar dado por la expresión:
      $$ f(x,y) = \ln(x^2 - y^2) $$
      Describa cualitativamente la geometría de la región del plano obtenida.
    }
    \pista{
      \textbf{Dominio:} El argumento del logaritmo debe ser estrictamente positivo ($x^2 - y^2 > 0$). Esto equivale a $x^2 > y^2 \implies |x| > |y|$. Geométricamente, representa el interior de los dos conos abiertos opuestos que contienen al eje $X$, delimitados por las rectas asíntotas $y = x$ e $y = -x$ (sin incluir las rectas). \\
      \textbf{Imagen:} Como la expresión $x^2 - y^2$ puede tomar cualquier valor dentro del intervalo $(0, \infty)$ bajo las condiciones del dominio, el logaritmo natural recorre todo su espectro. Por lo tanto, $\operatorname{im}(f) = \mathbb{R}$.
    }
  \end{ejerciciopropuesto}

  % --- 5. EJERCICIO PROPUESTO 3 \colon Raíz Cono Cerrado ---
  \begin{ejerciciopropuesto}{Regiones Cónicas y Fronteras Cerradas}{nivel-2}
    \enunciado{
      Considere el campo escalar:
      $$ f(x,y) = \sqrt{-x^2 + y^2} $$
      Halle analíticamente su dominio natural e imagen, y establezca la diferencia geométrica respecto al ejercicio anterior.
    }
    \pista{
      \textbf{Dominio:} El radicando exige que $y^2 - x^2 \geq 0 \implies y^2 \geq x^2$, lo cual se traduce en la desigualdad de valores absolutos $|y| \geq |x|$. Geométricamente, esto corresponde a las regiones (superior e inferior) que contienen al eje $Y$. A diferencia del problema anterior, esta región es \textit{cerrada}, lo que significa que sí incluye a las rectas fronteras $y = x$ e $y = -x$. \\
      \textbf{Imagen:} Al tratarse de una raíz cuadrada estándar cuyo radicando puede crecer indefinidamente conforme nos alejamos en el eje $Y$, el conjunto imagen corresponde a los reales no negativos: $\operatorname{im}(f) = [0, \infty)$.
    }
  \end{ejerciciopropuesto}

  % --- 6. EJERCICIO PROPUESTO 4 \colon Seno Acotado ---
  \begin{ejerciciopropuesto}{Periodicidad y Acotación Multivariable}{nivel-1}
    \enunciado{
      Para el campo escalar definido por la regla de correspondencia:
      $$ f(x,y) = \sin(x+y) $$
      Determine su dominio natural y su conjunto imagen.
    }
    \pista{
      \textbf{Dominio:} La función trigonométrica seno no impone ninguna restricción matemática sobre el comportamiento de sus argumentos, por lo que su dominio natural es todo el plano bidimensional, $\operatorname{dom}(f) = \mathbb{R}^2$. \\
      \textbf{Imagen:} Dado que la combinación lineal $x+y$ puede tomar cualquier valor real en el intervalo $(-\infty, \infty)$, y sabiendo que la función seno oscila de forma periódica, la imagen queda confinada de manera idéntica al caso unidimensional: $\operatorname{im}(f) = [-1, 1]$.
    }
  \end{ejerciciopropuesto}

  % --- 7. EJERCICIO PROPUESTO 5 \colon Raíz Parabólica ---
  \begin{ejerciciopropuesto}{Fronteras Parabólicas en el Plano}{nivel-2}
    \enunciado{
      Determine analíticamente el dominio natural y el conjunto imagen de la función:
      $$ f(x,y) = \sqrt{x^2 - y} $$
    }
    \pista{
      \textbf{Dominio:} La restricción de la raíz cuadrada de índice par exige que $x^2 - y \geq 0$, lo que equivale analíticamente a la inecuación $y \leq x^2$. En el plano cartesiano, esto representa geométricamente a todos los puntos que se encuentran sobre y por debajo de la parábola estándar $y = x^2$. \\
      \textbf{Imagen:} Debido a que la raíz cuadrada entrega exclusivamente valores no negativos y la diferencia $x^2 - y$ puede ser arbitrariamente grande (por ejemplo, fijando $x=0$ y haciendo que $y \to -\infty$), la imagen es $\operatorname{im}(f) = [0, \infty)$.
    }
  \end{ejerciciopropuesto}

\end{ejercicios}


% ==============================================================================
% PESTAÑA 5 \colon FÓRMULAS CLAVE (contentFormulas)
% ==============================================================================
\begin{formulas}

  \formula{Dominio de un Campo Escalar}
    {\operatorname{dom}(f) = \{ (x,y) \in \mathbb{R}^2 \colon f(x,y) \in \mathbb{R} \}}
    {El conjunto de todos los puntos en el plano para los cuales la regla de correspondencia de la función produce un valor real bien definido.}

  \formula{Imagen o Recorrido}
    {\operatorname{im}(f) = \{ z \in \mathbb{R} \colon z = f(x,y) \text{ para algún } (x,y) \in \operatorname{dom}(f) \}}
    {El conjunto de todos los valores numéricos (alturas, temperaturas, presiones) que la función efectivamente toma en el eje Z.}

  \formula{Restricciones de Dominio}
    {\begin{aligned}
      \text{Fracción: } \frac{1}{g} &\implies g \neq 0 \\[0.5em]
      \text{Logaritmo: } \ln(g) &\implies g > 0 \\[0.5em]
      \text{Raíz Par: } \sqrt{g} &\implies g \geq 0
    \end{aligned}}
    {Condiciones algebraicas obligatorias e indispensables para plantear el dominio natural: se debe recordar repasar hacia abajo las restricciones analíticas de denominadores no nulos, argumentos estrictamente positivos y radicandos no negativos.}

  \formula{Imágenes Elementales}
    {\begin{aligned}
      \operatorname{im}(e^u) &= (0, \infty) \\[0.5em]
      \operatorname{im}(\ln(u)) &= \mathbb{R} \\[0.5em]
      \operatorname{im}(\sqrt{u}) &= [0, \infty)
    \end{aligned}}
    {Comportamiento analítico y recorridos canónicos verticales de las funciones base de una variable. Recordar de memoria estos intervalos es clave para construir la imagen final de campos escalares complejos.}

  \formula{Frontera Circular Típica}
    {x^2 + y^2 = r^2}
    {Ecuación de la circunferencia de radio $r$ centrada en el origen, la cual suele aparecer como la frontera geométrica al despejar restricciones de raíces o logaritmos radiales.}

\end{formulas}

\end{document}